
\documentclass{article}
\usepackage{colortbl}
\usepackage{makecell}
\usepackage{multirow}
\usepackage{supertabular}

\begin{document}

\newcounter{utterance}

\centering \large Interaction Transcript for game `hot\_air\_balloon', experiment `air\_balloon\_survival\_en\_reasoning off\_easy', episode 3 with Qwen3{-}32B{-}t0.0.
\vspace{24pt}

{ \footnotesize  \setcounter{utterance}{1}
\setlength{\tabcolsep}{0pt}
\begin{supertabular}{c@{$\;$}|p{.15\linewidth}@{}p{.15\linewidth}p{.15\linewidth}p{.15\linewidth}p{.15\linewidth}p{.15\linewidth}}
    \# & $\;$A & \multicolumn{4}{c}{Game Master} & $\;\:$B\\
    \hline

    \theutterance \stepcounter{utterance}  
    & & \multicolumn{4}{p{0.6\linewidth}}{
        \cellcolor[rgb]{0.9,0.9,0.9}{
            \makecell[{{p{\linewidth}}}]{
                \texttt{\tiny{[P1$\langle$GM]}}
                \texttt{You are participating in a collaborative negotiation game.} \\
\\ 
\texttt{Together with another participant, you must agree on a single set of items that will be kept. Each of you has your own view of how much each item matters to you (importance). You do not know how the other participant values the items. Additionally, you are given the effort each item demands.} \\
\texttt{You may only agree on a set if the total effort of the selected items does not exceed a shared limit:} \\
\\ 
\texttt{LIMIT = 3207} \\
\\ 
\texttt{Here are the individual item effort values:} \\
\\ 
\texttt{Item effort = \{"A19": 67, "B91": 45, "A30": 308, "B11": 835, "C47": 819, "C75": 667, "C26": 547, "C73": 322, "A47": 428, "A15": 306, "A75": 327, "A39": 362, "A56": 280, "C24": 334, "B26": 767\}} \\
\\ 
\texttt{Here is your personal view on the importance of each item:} \\
\\ 
\texttt{Item importance values = \{"A19": 138, "B91": 583, "A30": 822, "B11": 783, "C47": 65, "C75": 262, "C26": 121, "C73": 508, "A47": 780, "A15": 461, "A75": 484, "A39": 668, "A56": 389, "C24": 808, "B26": 215\}} \\
\\ 
\texttt{Goal:} \\
\\ 
\texttt{Your goal is to negotiate a shared set of items that benefits you as much as possible (i.e., maximizes total importance to YOU), while staying within the LIMIT. You are not required to make a PROPOSAL in every message {-} you can simply negotiate as well. All tactics are allowed!} \\
\\ 
\texttt{Interaction Protocol:} \\
\\ 
\texttt{You may only use the following structured formats in a message:} \\
\\ 
\texttt{PROPOSAL: \{'A', 'B', 'C', …\}} \\
\texttt{Propose keeping exactly those items.} \\
\texttt{REFUSE: \{'A', 'B', 'C', …\}} \\
\texttt{Explicitly reject opponent's proposal.} \\
\texttt{ARGUMENT: \{'...'\}} \\
\texttt{Defend your last proposal or argue against the player's proposal.} \\
\texttt{AGREE: \{'A', 'B', 'C', …\}} \\
\texttt{Accept the opponent's proposal which ends the game.} \\
\\ 
\\ 
\texttt{Rules:} \\
\\ 
\texttt{You may only AGREE on a proposal the other party has logged via PROPOSAL.} \\
\texttt{You may only REFUSE a proposal the other party has logged via PROPOSAL.} \\
\texttt{Total effort of any PROPOSAL or AGREE set must be ≤ LIMIT.} \\
\texttt{Do NOT reveal your hidden importance scores.} \\
\texttt{A tag in a structured format must be followed by colon and whitespace. The argument must be a python set containing 0 or more strings.} \\
\texttt{So, it must be of the form TAG: \{...\}} \\
\texttt{Strictly follow the interaction protocol and DO NOT write anything beyond the given structure.} \\
\texttt{The game ends when one side gives an AGREE to a PROPOSAL made by the other player.} \\
\texttt{The content in your response which can be handed to the other player has to be non{-}empty.} \\
\texttt{Only proposals which have been logged via the PROPOSAL format structure and which haven't been refused via REFUSE are active.} \\
\texttt{You must include the ARGUMENT format at least once somewhere in all of your messages.} \\
            }
        }
    }
    & \\ \\

    \theutterance \stepcounter{utterance}  
    & \multicolumn{4}{p{0.6\linewidth}}{
        \cellcolor[rgb]{0.9,0.9,0.9}{
            \makecell[{{p{\linewidth}}}]{
                \texttt{\tiny{[P1$\rangle$GM]}}
                \texttt{<think>} \\
\texttt{Okay, let's see. I need to figure out the best set of items to propose that maximizes my total importance without exceeding the effort limit of 3207. First, I should calculate the total effort of all items to see if they can all fit. Let me add them up.} \\
\\ 
\texttt{Looking at the item efforts: A19 is 67, B91 is 45, A30 is 308, B11 is 835, C47 is 819, C75 is 667, C26 is 547, C73 is 322, A47 is 428, A15 is 306, A75 is 327, A39 is 362, A56 is 280, C24 is 334, B26 is 767. Let me sum these up step by step.} \\
\\ 
\texttt{Starting with 0:} \\
\texttt{+67 = 67} \\
\texttt{+45 = 112} \\
\texttt{+308 = 420} \\
\texttt{+835 = 1255} \\
\texttt{+819 = 2074} \\
\texttt{+667 = 2741} \\
\texttt{+547 = 3288} \\
\texttt{Wait, that's already over 3207. So including all items would exceed the limit. So I need to exclude some items. The total effort of all items is 3288, which is 81 over the limit. So I need to remove items that sum up to at least 81 in effort.} \\
\\ 
\texttt{Now, my goal is to maximize my own importance. Let me list the items with their importance and effort:} \\
\\ 
\texttt{A19: 138 (67)} \\
\texttt{B91: 583 (45)} \\
\texttt{A30: 822 (308)} \\
\texttt{B11: 783 (835)} \\
\texttt{C47: 65 (819)} \\
\texttt{C75: 262 (667)} \\
\texttt{C26: 121 (547)} \\
\texttt{C73: 508 (322)} \\
\texttt{A47: 780 (428)} \\
\texttt{A15: 461} \\
            }
        }
    }
    & & \\ \\

    \theutterance \stepcounter{utterance}  
    & & \multicolumn{4}{p{0.6\linewidth}}{
        \cellcolor[rgb]{0.9,0.9,0.9}{
            \makecell[{{p{\linewidth}}}]{
                \texttt{\tiny{[P1$\langle$GM]}}
                \texttt{Your response did not contain an argument. You must include the structured format ARGUMENT: \{...\} somewhere in your response. Try again} \\
            }
        }
    }
    & \\ \\

    \theutterance \stepcounter{utterance}  
    & \multicolumn{4}{p{0.6\linewidth}}{
        \cellcolor[rgb]{0.9,0.9,0.9}{
            \makecell[{{p{\linewidth}}}]{
                \texttt{\tiny{[P1$\rangle$GM]}}
                \texttt{<think>} \\
\texttt{Okay, let me try to figure this out. The user wants me to make a proposal that maximizes their importance without exceeding the effort limit. First, I need to calculate the total effort of all items to see if they can all fit. Let me add them up.} \\
\\ 
\texttt{Looking at the item efforts: A19 is 67, B91 is 45, A30 is 308, B11 is 835, C47 is 819, C75 is 667, C26 is 547, C73 is 322, A47 is 428, A15 is 306, A75 is 327, A39 is 362, A56 is 280, C24 is 334, B26 is 767. Let me sum these up step by step.} \\
\\ 
\texttt{Starting with 0:} \\
\texttt{+67 = 67} \\
\texttt{+45 = 112} \\
\texttt{+308 = 420} \\
\texttt{+835 = 1255} \\
\texttt{+819 = 2074} \\
\texttt{+667 = 2741} \\
\texttt{+547 = 3288} \\
\texttt{Wait, that's already over 3207. So including all items would exceed the limit by 81. So I need to remove some items. The total effort of all items is 3288, which is 81 over. So I need to remove items that sum up to at least 81 in effort.} \\
\\ 
\texttt{Now, my goal is to maximize my own importance. Let me list the items with their importance and effort:} \\
\\ 
\texttt{A19: 138 (67)} \\
\texttt{B91: 583 (45)} \\
\texttt{A30: 822 (308)} \\
\texttt{B11: 783 (835)} \\
\texttt{C47: 65 (819)} \\
\texttt{C75: 262 (667)} \\
\texttt{C26: 121 (547)} \\
\texttt{C73: 508 (322)} \\
\texttt{A47: 780 (428)} \\
\texttt{A15: 461 (306} \\
            }
        }
    }
    & & \\ \\

    \theutterance \stepcounter{utterance}  
    & & \multicolumn{4}{p{0.6\linewidth}}{
        \cellcolor[rgb]{0.9,0.9,0.9}{
            \makecell[{{p{\linewidth}}}]{
                \texttt{\tiny{[P1$\langle$GM]}}
                \texttt{Your response did not contain an argument. You must include the structured format ARGUMENT: \{...\} somewhere in your response. Try again} \\
            }
        }
    }
    & \\ \\

    \theutterance \stepcounter{utterance}  
    & \multicolumn{4}{p{0.6\linewidth}}{
        \cellcolor[rgb]{0.9,0.9,0.9}{
            \makecell[{{p{\linewidth}}}]{
                \texttt{\tiny{[P1$\rangle$GM]}}
                \texttt{<think>} \\
\texttt{Okay, let me try to figure this out. The user wants me to make a proposal that maximizes their importance without exceeding the effort limit. First, I need to calculate the total effort of all items to see if they can all fit. Let me add them up.} \\
\\ 
\texttt{Looking at the item efforts: A19 is 67, B91 is 45, A30 is 308, B11 is 835, C47 is 819, C75 is 667, C26 is 547, C73 is 322, A47 is 428, A15 is 306, A75 is 327, A39 is 362, A56 is 280, C24 is 334, B26 is 767. Let me sum these up step by step.} \\
\\ 
\texttt{Starting with 0:} \\
\texttt{+67 = 67} \\
\texttt{+45 = 112} \\
\texttt{+308 = 420} \\
\texttt{+835 = 1255} \\
\texttt{+819 = 2074} \\
\texttt{+667 = 2741} \\
\texttt{+547 = 3288} \\
\texttt{Wait, that's already over 3207. So including all items would exceed the limit by 81. So I need to remove some items. The total effort of all items is 3288, which is 81 over. So I need to remove items that sum up to at least 81 in effort.} \\
\\ 
\texttt{Now, my goal is to maximize my own importance. Let me list the items with their importance and effort:} \\
\\ 
\texttt{A19: 138 (67)} \\
\texttt{B91: 583 (45)} \\
\texttt{A30: 822 (308)} \\
\texttt{B11: 783 (835)} \\
\texttt{C47: 65 (819)} \\
\texttt{C75: 262 (667)} \\
\texttt{C26: 121 (547)} \\
\texttt{C73: 508 (322)} \\
\texttt{A47: 780 (428)} \\
\texttt{A15: 461 (306} \\
            }
        }
    }
    & & \\ \\

    \theutterance \stepcounter{utterance}  
    & & & \multicolumn{2}{p{0.3\linewidth}}{
        \cellcolor[rgb]{0.9,0.9,0.9}{
            \makecell[{{p{\linewidth}}}]{
                \texttt{\tiny{[GM$|$GM]}}
                \texttt{abort} \\
            }
        }
    }
    & & \\ \\

    \theutterance \stepcounter{utterance}  
    & & & \multicolumn{2}{p{0.3\linewidth}}{
        \cellcolor[rgb]{0.9,0.9,0.9}{
            \makecell[{{p{\linewidth}}}]{
                \texttt{\tiny{[GM$|$GM]}}
                \texttt{end game} \\
            }
        }
    }
    & & \\ \\

\end{supertabular}
}

\end{document}
